\section{总结与延伸}

Lecture 2 建立了大模型训练的最小资源语言：

\begin{enumerate}
    \item Tensor 是 data、parameters、gradients、optimizer state、activations 的统一载体。
    \item Memory 由 numel、dtype、device 和生命周期决定。
    \item FLOPs 是计算量，FLOP/s 是速度，MFU 衡量有效硬件利用率。
    \item Arithmetic intensity 判断 memory-bound 和 compute-bound。
    \item 训练一步常用 \(6BP\) 粗估，但 memory 还要算 optimizer state 和 activations。
    \item Gradient accumulation、activation checkpointing、ZeRO/FSDP 都是资源交换。
\end{enumerate}

\begin{importantbox}{最终 takeaway}
Resource accounting 是 CS336 的工程底座。不会算 FLOPs、memory 和 bandwidth，就无法判断一个训练 recipe 是否可行，也无法理解后续 kernels、parallelism、inference 和 scaling laws 为什么这样设计。
\end{importantbox}

\subsection{拓展阅读}

\begin{itemize}
    \item JAX Scaling Book: roofline and systems chapters.
    \item NVIDIA H100/B200 datasheets and Transformer Engine FP8 documentation.
    \item PyTorch AMP, activation checkpointing, and optimizer documentation.
    \item Transformer memory and FLOPs accounting blog posts.
\end{itemize}

\subsection{Executable source 索引：原始教学节点如何落到讲义}

下面的索引保留 executable lecture 中会被用作导航标题的原始名称。左列不是额外知识点清单，而是把源码节点与前文教学展开一一对应，便于读者回到 \texttt{lecture02-slides.py} 复核上下文、公式和代码。

\begin{longtable}{p{0.36\textwidth}p{0.54\textwidth}}
\toprule
\textbf{Executable-source title} & \textbf{本讲中的教学处理} \\
\midrule
\texttt{Mixed precision} & 数值格式一节：bf16 参数/激活/梯度、fp32 optimizer states，以及 AMP 对 matmul 与数值敏感算子的区别处理。 \\
\texttt{Intuitions} & Compute Accounting：从训练规模数量级进入 FLOPs、FLOP/s 与可达性能直觉。 \\
\texttt{Linear model} & 矩阵乘法与梯度例子：从 \(Y=XW\) 和两层线性网络逐项计算 forward/backward FLOPs。 \\
\texttt{Model FLOPs utilization (MFU)} & 实际 FLOP/s 与硬件承诺峰值的比值、约 0.5 的经验阈值及其忽略通信的边界。 \\
\texttt{Zoom in on one layer} & 放大 \(h_2=h_1W_2\)，分别计算 activation gradient 与 weight gradient。 \\
\texttt{Consider all layers} & 把单层结论推广到所有参数，得到训练一步约 \(6BP\) FLOPs。 \\
\texttt{Compute (for one training step)} & Optimizer state 与 training loop：区分 forward、backward 和 parameter update 的计算与状态读写。 \\
\texttt{Transformers} & 说明 \(6BP\) 对短上下文 Transformer 是常用近似，同时提醒长上下文 attention 需要补充序列长度相关项。 \\
\bottomrule
\end{longtable}

\begin{knowledgebox}{如何使用这张 source-node 索引}
若目标是学习资源核算，按正文顺序阅读即可；若目标是验证覆盖或继续扩写，则用左列精确标题搜索 executable source，再检查标题附近的 \texttt{text(...)}, tensor 操作和公式。索引只负责导航，真正的解释仍在前文的推导、读图框、代码和 teacher-voice 段落中。
\end{knowledgebox}

\subsection{自测：能否独立完成四本资源账}

读完本讲后，不应只会复述术语，而应能在没有 profiler 的情况下先写出可检查的数量级预测。下面四题分别检验 compute、memory、bandwidth 与 tradeoff；建议先遮住右列，独立写出公式、假设和遗漏项。

\begin{longtable}{p{0.42\textwidth}p{0.48\textwidth}}
\toprule
\textbf{自测问题} & \textbf{答案线索与验收点} \\
\midrule
已知参数量 \(P\)、token 数 \(D\)、GPU 数 \(G\)、单卡峰值 \(R\) 与 MFU，如何估算训练天数？ & 先算 \(6PD\)，再除以 \(G\times R\times\mathrm{MFU}\)；必须写明 dtype 峰值口径，并提醒通信、checkpoint 与故障会拉长 wall-clock。 \\
一张卡能否放下某个 AdamW 训练配置？ & 至少分开列 parameters、gradients、\(m\)、\(v\)、activations 与 temporary buffers；不能用“参数量乘 2 bytes”代替完整训练显存。 \\
为什么大 matmul 可能 compute-bound，而逐元素激活通常 memory-bound？ & 比较 arithmetic intensity 与 accelerator intensity；说明 matmul 的输入复用如何提高每 byte FLOPs，以及 precision 会同时改变两边的比值。 \\
显存不足时应先选 gradient accumulation、activation checkpointing 还是 ZeRO/FSDP？ & 先确认峰值来自 activations 还是复制的训练状态；三者分别主要用额外 micro-steps、重算或通信换显存，选择必须回到端到端吞吐。 \\
\bottomrule
\end{longtable}

\begin{warningbox}{自测的最低合格标准}
每个答案都必须同时写出公式、单位、假设和未计入项。只有一个数字而没有口径，不能算 resource accounting；只有工具实测而没有事前预测，也无法判断结果是否合理。
\end{warningbox}

\end{document}
